% ************************** Thesis Abstract *****************************
% The PhD school of Politecnico di Torino limits the abstract to 4000 characters.
\begin{abstract}
Innovative reactor concepts such as lead-cooled fast reactors and heat-pipe
micro reactors are not served, as it stands, by the computational chain
consolidated on light-water reactors, and the Monte Carlo solution of the full
core remains their reference, at a cost that the many states required by design
and operation cannot afford. This thesis works on the data that every practical
scheme places between that reference and the many evaluations: the multi-group
constants read by a deterministic solver, and a database of fission matrices
from which any state of the core is recovered by interpolation. In both cases
the data are prepared once from a few reference calculations, corrected without
modifying the solver that uses them, and every result is verified against a
Monte Carlo calculation made for that purpose.

The first Part concerns the multi-group data of lead-cooled fast reactors. For a
simplified model of the Westinghouse LFR the constants are generated with
Serpent~2 along two routes, a single full-core model and a set of mini-core
models representative of the regions of the system, and used in a
finite-element diffusion solver with a superhomogenization correction of the
regions responsible for the largest discrepancies. The mini-core route proves
the better choice, for accuracy and for computing time alike, whenever several
control-rod configurations have to be analysed, and the correction brings the
multiplication factor and the power distribution close to the reference. For
the ALFRED core the energy grid on which the constants are condensed is treated
as the unknown of an optimization. A neural metamodel trained on $91\,723$
full-core evaluations of the FRENETIC code replaces the solver inside a genetic
search, making one thousand independent searches affordable where three had
been, and every structure the searches return is recomputed with the solver.
The dependence of the result on the random seed is measured, the best
fifteen-group structure improves on the three previously published ones, and
the optimized structures share an exclusive choice in the fast region and place
their boundaries according to the neutron importance rather than to the flux.
Their quality rests on a cancellation between the condensation error and the
error of the deterministic model, which binds the result to the core, the
solver and the objective, while the procedure carries over.

The second Part develops fission matrix methods for the online monitoring of a
heat-pipe micro reactor, on the model distributed on the Virtual Test Bed. A
database of matrices over fuel temperature and control-drum angle, with a
prescribed correlation between power and temperature, supports a criticality
search that converges in about one second on a single core and returns states
whose multiplication factor lies within a few tens of pcm of a reference Monte
Carlo calculation; a multiplicative burnup correction ratio extends it along
the fuel cycle without adding an axis to the database. Two further correction
ratios account for the temperature difference between core and reflector and
for sharp temperature gradients inside the core, and the statistical
uncertainty that the interpolated matrices inherit from the tallies is derived
in closed form and validated against independent replicas of the whole
database.
\end{abstract}
